\question
\tocsection{Designing a Dam}

How might you design a dam?

\begin{hints}
	\begin{enumerate}
		\item What materials should be used?
		\item What shape should the dam be? Describe the cross-section.
	\end{enumerate}
\end{hints}

\begin{solution}
This question was deliberately open-ended so interviewees were expected to come up with what aspects would be important to consider. The following are suggested.
\begin{itemize}
	\item Porosity. The material would need to be non-porous so that water does not leak through. This rules out cheap materials such as sand and earth.
	\item Corrosion. Water can react with many materials which can increase maintenance requirements or compromise the dam integrity.
	\item Cost. Dams are very large so making them out of exotic materials or with an intricate design would be expensive. Even materials such as steel would be prohibitively expensive given the quantities that are needed.
	\item Strength. The dam will need to resist a lot of force.
\end{itemize}

Candidates are anticipated to arrive at a giant uniform mass of concrete or other suitable material. It remains to determine what shape the dam should be so that materials are used efficiently. Water has a high density which means the pressure that a dam needs to hold back can be very large - understanding this pressure is what determines the shape of the dam.

Stevin's law states that the rate of change of pressure in the vertical direction within a static fluid is directly proportional to the specific weight of the fluid. As water can be approximated as incompressible we can assume a constant density, therefore the hydrostatic pressure increases linearly with depth of water. Pascal's law states that the pressure in a fluid acts equally in all directions, therefore the pressure would also act against the walls of the dam.

For the dam to withstand the higher pressure it would need to be stronger towards the base. We can assume that the strength of a wall (that is, how much pressure it can withstand) is proportional to its thickness, therefore the dam wall thickness will increase linearly with the depth of the water.

\end{solution}